\section{Hamiltonians and operators in a nonorthogonal basis}
\label{sec:operators}
\subsection{Covariant matrices and the physical metric}
Let $\Phi$ denote the row of basis functions $|\phi_\mu\rangle$, and
$|\psi\rangle=\Phi c$. Define
\begin{equation}
 S=\Phi^\dagger\Phi,\qquad
 H=\Phi^\dagger\widehat H\Phi,\qquad
 X_a=\Phi^\dagger\widehat r_a\Phi .
 \label{eq:ao-data}
\end{equation}
The overlap must be positive definite on the retained basis, and the
normalization is $c^\dagger Sc=1$. The generalized eigenproblem is
\begin{equation}
 HC=SC\varepsilon,\qquad C^\dagger SC=I .
 \label{eq:generalized}
\end{equation}
A covariant matrix $A$ acts on coefficient vectors through $S^{-1}A$.
Consequently, composition of two \emph{projected} operators is represented
by $AS^{-1}B$, and the projected commutator by
$AS^{-1}B-BS^{-1}A$. An ordinary product $AB$ is generally incorrect.

For a nonsingular change of basis $\Phi'=\Phi T$,
\begin{equation}
 S'=T^\dagger ST,\quad A'=T^\dagger AT,\quad c'=T^{-1}c .
 \label{eq:basis-change}
\end{equation}
The expectation $c^\dagger Ac$ and the generalized spectrum are unchanged.
The orthonormal representation
$\widetilde A=S^{-1/2}AS^{-1/2}$ is useful for proofs and dense reference
calculations. It need not be formed in a sparse implementation: products
can be evaluated by solving linear systems with $S$.

A finite-basis projection is not an algebra homomorphism for arbitrary
continuum operators. With $\Pi=\Phi S^{-1}\Phi^\dagger$,
$\Phi^\dagger\widehat A\Pi\widehat B\Phi$ differs in general from
$\Phi^\dagger\widehat A\widehat B\Phi$.
Thus a commutator or a variance of projected positions is an observable
of the projected model, not automatically its complete-basis continuum
counterpart. This distinction is monitored through basis convergence
rather than hidden in a numerical tolerance.

\subsection{Finite positions and AO export}
For an isolated system, analytic Gaussian first moments provide
$X_a$ in the same basis ordering as $H$ and $S$. Translation of the
coordinate origin by $x_a$ gives $X_a-x_aS$, not $X_a-x_aI$.
Likewise an energy query uses $H-ES$. These two simple identities are
essential for origin and energy-shift invariance.

The \texttt{AO\_MATRICES} print section can export overlap, scalar KS
Hamiltonian, and the three position matrices. They permit an independent
finite-Hamiltonian calculation. Activating SOC in a property section
does not silently turn the separate scalar KS export into a spinor
Hamiltonian export. Model Hamiltonians can use the low-level matrix
kernels when the metric, positions, and symmetry structure are supplied
consistently; this does not imply automatic adapters for every
electronic-structure method.

\subsection{Bloch operators and directed overlaps}
In the cell Fourier convention,
\begin{equation}
 H_{\mu\nu}(\bk)=\sum_{\bR}e^{\ii\bk\cdot\bR}H_{\mu\nu}(\bR),
 \qquad
 S_{\mu\nu}(\bk)=\sum_{\bR}e^{\ii\bk\cdot\bR}S_{\mu\nu}(\bR).
 \label{eq:bloch-matrices}
\end{equation}
Here $\bk$ is Cartesian and includes inverse-length units; fractional
reciprocal coordinates instead require $2\pi$ in the phase.
Image-cell labels cannot be discarded before forming derivatives or
inter-k-point connections.

For cell-periodic states and a directed step $\bb$, the occupied-state
overlap is
\begin{equation}
 M(\bk,\bb)=C^\dagger(\bk)O(\bk+\bb;\bb)C(\bk+\bb),\qquad
 O_{\mu\nu}=\sum_{\bR}e^{\ii(\bk+\bb)\cdot\bR}
 \langle\phi_{\mu0}|e^{-\ii\bb\cdot\br}|\phi_{\nu\bR}\rangle .
 \label{eq:wannier90-mmn-definition}
\end{equation}
The zero-step limit is $S(\bk)$, but a finite link is not obtained by
replacing $O$ with $S$, or by simply contracting coefficient vectors.
Neighbor lists, reciprocal-boundary phases, and the selected band rank
are part of the observable. Wannier90's file and library interfaces can
consume such data \cite{Pizzi2020}; native Wilson calculations need no
Wannier localization or interpolation.

\subsection{Spin--orbit coupling and symmetry}
The implementations considered here start from a converged scalar
restricted SCF potential and add GTH spin--orbit matrix elements.
The localizer and Bloch-Kubo paths use the complete AO spinor space,
whereas a second-variational band-topology calculation uses a selected
scalar-state space whose size must be converged separately.
Neither construction is self-consistent noncollinear SOC DFT.

In spin-major ordering, the overlap is
$\mathbb S=\operatorname{diag}(S,S)$ and physical time reversal is
$\Theta=JK$, where
\begin{equation}
 J=\begin{pmatrix}0&I\\-I&0\end{pmatrix},\qquad
 \Theta^2=-I .
\end{equation}
Time-reversal relations must be checked for the Hamiltonian, metric,
and relevant operators; the presence of an SOC pseudopotential alone
does not establish them.
A Cartesian spatial rotation $R$ acts on spin through an SU(2) lift of
$\det(R)R$, with inversion acting trivially on spin and separately on
orbital parity. The sign ambiguity of this lift is retained in the
double-group multiplication law.
